\section{Theory}

\subsection{The Bio-Gravitational Number and the Allometric Trap}

For a vertebral column of length $L$, flexural stiffness $EI$, mass $M$, and gravitational acceleration $g$, we define the \textbf{Bio-Gravitational Number}:
\begin{equation}
\mathcal{B}_g = \frac{EI}{M g L^2},
\label{eq:bg_number_passive}
\end{equation}
a dimensionless ratio of passive flexural resistance to gravitational bending moment. When $\mathcal{B}_g > 0.1$, an organism can resist gravitational collapse through passive stiffness alone. When $\mathcal{B}_g < 0.1$, the organism enters the ``Allometric Trap'': it must actively maintain its shape at continuous metabolic cost. Cross-species analysis reveals that small quadrupeds operate well above this threshold, while humans ($\mathcal{B}_g \approx 0.01$) are deep within the trap, relying almost entirely on active metabolic maintenance to prevent spinal collapse~\cite{smit2020adolescent}. The Bio-Gravitational Number is analogous to the Froude number for locomotion, defining a fundamental physical boundary for spinal stability. However, scaling exponents vary systematically depending on physiological state~\cite{glazier2022variable}, and vertebral scaling exhibits size-dependent breakpoints~\cite{smith2024small}, emphasizing that bipedalism and human-specific axial loading uniquely exacerbate this vulnerability~\cite{gorman2009idiopathic}.

The central problem of vertebrate morphogenesis is the translation of a one-dimensional genetic code into a robust three-dimensional geometry that can withstand environmental loads. Classical structural mechanics predicts that a flexible column clamped at the base and subjected to its own weight will buckle or sag into a monotonic C-shape~\cite{goriely2017mathematics}. Yet, the healthy human spine exhibits a complex, alternating S-shape (lordosis and kyphosis).

This anomaly suggests that the S-shape is not a passive equilibrium but an active \emph{Counter-Curvature}---a geometric standing wave maintained by the continuous expenditure of metabolic energy against the gravitational field. This maintenance constitutes a biological implementation of \emph{optimal feedback control}~\cite{todorov2002optimal} and can be conceptually modeled using \emph{Latent Imagination}~\cite{hafner2019dreamer}, where the nervous system learns an internal world model of gravity and biomechanics to predict and optimize long-horizon postural behaviors by minimizing a cost function comprising both ``geometric error'' (deviation from the target shape) and ``metabolic effort'' (ATP consumption).

\subsection{The Recovery Ratchet: Geometric Restoration versus Growth-Remodeling}

The onset of AIS is best understood not as a balance failure but as a \emph{geometry-recovery} failure. Developmental growth continually perturbs the spine from its reference shape; the active system must restore that shape faster than growth cements the deviation. We separate the standing curvature deviation into a recoverable elastic component $\kappa_e$ --- the part that reduces when gravitational load is removed (supine, or on bending films) --- and a permanently remodeled component $\kappa_p$. Two rates compete: the active geometric-restoration rate $k_r$, and the growth-remodeling rate $k_g(t)$, by which a sustained deviation is cemented through Hueter--Volkmann growth-plate modulation~\cite{stokes2006hueter, villemure2009growth}. For a perturbation that recoils at rate $k_r$ and is cemented at rate $k_g$, the fraction converted to permanent set is exactly $k_g/(k_r+k_g)$, giving
\begin{equation}
\frac{d\kappa_p}{dt} = k_g(t)\,\kappa_e\,\frac{k_g(t)}{k_r + k_g(t)}, \qquad \mathrm{flexibility}(t) = \frac{\kappa_e}{\kappa_e + \kappa_p(t)}.
\label{eq:recovery_ratchet}
\end{equation}
When $k_r \gg k_g$ the spine recoils and almost nothing sets ($\kappa_p \to 0$, flexibility $\to 1$). When $k_g \gtrsim k_r$ --- as occurs when $k_g(t)$ peaks at peak height velocity (PHV) --- each perturbation is partly cemented before it can be recovered, and the deviation \emph{ratchets} monotonically into a permanent curve (Fig.~\ref{fig:recovery_ratchet}). Because $k_g(t)$ is maximal precisely during the growth spurt, the ratchet predicts the clustering of AIS onset at PHV \emph{without} invoking any postural-balance deficit --- consistent with the clinical observation that affected adolescents are otherwise neurologically normal and curves are detected incidentally. The model yields a falsifiable structural signature: \emph{curve flexibility (reducibility) declines as $\kappa_p$ accumulates}, so progressing curves should become progressively less reducible on supine and bending films, and low reducibility at presentation should predict progression. The genetic risk factors associated with AIS (reduced collagen-dependent damping, altered mechanosensation) act in this framework by lowering $k_r$, shifting the individual toward the ratchet regime.

\subsection{Complementary Fast-Loop Delay Dynamics}

The slow recovery ratchet is complemented by a faster postural-control loop. For completeness we retain a delay-differential treatment of that loop --- proprioceptive proportional-derivative (PD) feedback subject to neural delay $\tau$~\cite{milton2009time, insperger2013acceleration} --- developed fully in the Supplementary Material; it is summarized here only to position it as a secondary, intact loop rather than the primary onset mechanism. We formalize the fast loop as a Delay Differential Equation (DDE).

The proprioceptive control law generating corrective force $F_{control}$ with feedback delay $\tau$ is:
\begin{equation}
\boldsymbol{\kappa}_{\mathrm{rest}}(s) = \boldsymbol{\kappa}_{\mathrm{gen}} + \chi_\kappa \nabla I(s).
\label{eq:kappa_rest_field}
\end{equation}
where $K_P$ and $K_D$ are proportional and derivative gains, and the error is $e(s,t) = \kappa_{measured}(s,t) - \kappa_{target}$.

Linearizing around the straight configuration, the dynamics of the spinal column become:
\begin{equation}
\rho A \frac{\partial^2 \kappa}{\partial t^2} + \eta \frac{\partial \kappa}{\partial t} + EI \frac{\partial^4 \kappa}{\partial s^4} + K_P \kappa(s, t-\tau) + K_D \frac{\partial \kappa}{\partial t}(s, t-\tau) = 0
\label{eq:dde_spine}
\end{equation}

The stability of this active structure is governed by the dimensionless Bio-Gravitational Number $\mathcal{B}_g$:
\begin{equation}
\mathcal{B}_g = \frac{\chi_M \langle |\nabla I| \rangle}{\rho A g L^2}.
\label{eq:bg_number_info}
\end{equation}
Crucially, when $\mathcal{B}_g > 1$, the biological system effectively dominates gravitational mechanics. At the molecular level, this coupling relies on the structural rigidity of patterning transcription factors; AlphaFold structural analysis of critical HOX proteins (e.g., HOXC8, UniProt P31273; HOXB13, UniProt Q92826) reveals conserved, high-confidence DNA-binding domains capable of enforcing the sharp spatial identity boundaries required by our field parameterization.

During the adolescent growth spurt, the spine elongates rapidly, and the neural feedback delay $\tau$ rises because nerve conduction velocity does not scale proportionately with height~\cite{lang1985nerve, robinson2001nerve}. In the fast loop this manifests as a \textbf{Derivative Gain Trap}: derivative gains ($K_D$) that provide damping in childhood lose phase margin as $\tau$ grows. We stress that, on its own, this fast-loop margin loss would produce \emph{oscillatory} instability (postural sway) --- the signature of global sensorimotor disease (e.g.\ cerebral palsy), which AIS patients do not exhibit. We therefore treat the fast loop as a modulator of the slow recovery ratchet rather than the onset mechanism: rising $\tau$ and lost damping degrade the effective restoration rate $k_r$, tipping the ratchet of Eq.~\ref{eq:recovery_ratchet} toward progressive cementing.

\subsection{Coupling to the Information-Elastic Continuum}

The deviation that the recovery ratchet acts upon is set by the underlying information-elastic geometry. Below we specify how developmental patterning enters the rod's rest curvature and stiffness.

Assuming solutions of the form $\kappa(s,t) = A e^{\lambda t} \sin\left(\frac{n\pi s}{L}\right)$, the characteristic equation is:
\begin{equation}
\boldsymbol{\kappa}_{\mathrm{rest}}(s) = \boldsymbol{\kappa}_{\mathrm{gen}} + \chi_\kappa \nabla I(s),
\label{eq:kappa_rest_iec}
\end{equation}
where $\chi_\kappa$ is a geometric coupling constant and $\boldsymbol{\kappa}_{\mathrm{gen}}$ is the baseline genetic curvature.

\paragraph{Stiffness modulation (IEC-2).}
The information field modulates the effective bending and torsional stiffness:
\begin{equation}
B_{\mathrm{eff}}(s) = E_0 I_{\mathrm{area}} (1 + \chi_E I(s)), \qquad
C_{\mathrm{eff}}(s) = G_0 J (1 + \chi_C I(s)).
\label{eq:stiffness_modulation}
\end{equation}

The total energy functional governing the rod's equilibrium is:
\begin{equation}
\mathcal{E}_{\mathrm{total}} = \int_0^L \left[ \frac{1}{2} B_{\mathrm{eff}}(\kappa - \kappa_{\mathrm{rest}})^2 + \frac{1}{2} C_{\mathrm{eff}} \tau^2 + \frac{1}{2} K_{\mathrm{eff}} \varepsilon^2 + \rho A \mathbf{r} \cdot \mathbf{g} \right] ds.
\label{eq:iec_energy}
\end{equation}

For purely imaginary roots $\lambda = \pm i\omega$ the fast loop reaches its damping margin. The reason AIS coincides with the growth spurt is, however, the recovery ratchet of Eq.~\ref{eq:recovery_ratchet}: the growth-remodeling rate $k_g(t)$ peaks at PHV and overtakes the geometric-restoration rate $k_r$, so transient deviation is cemented faster than it can be recovered precisely in the developmental window where AIS initiates.

\subsection{The Energy Deficit Window: NAD+ and Proprioceptive Integrity}

The biological spine is a dissipative structure requiring continuous metabolic expenditure to maintain its sensory and structural integrity. The metabolic power required to sustain this configuration scales superlinearly with length ($L^4$ demand vs $L^2$ surface-limited supply)~\cite{sato2025convective}.

This scaling mismatch creates a transient \textbf{Energy Deficit Window} during the adolescent growth spurt. We propose that this metabolic deficit directly lowers the geometric-restoration rate $k_r$ of the recovery ratchet (Eq.~\ref{eq:recovery_ratchet}). Specifically, proprioceptive Ia/II afferents and the paraspinal actuators that effect restoration are among the largest and most metabolically demanding tissues.

Recent evidence demonstrates that NAD+ depletion during development causes severe spinal deformities~\cite{basse2021nampt}, while oxidative stress directly drives scoliosis in animal models~\cite{pumputis2025oxidative}. Furthermore, NAD+ insufficiency activates the SARM1 pathway, which triggers local axon degeneration~\cite{gerdts2015sarm1}. We hypothesize that the adolescent Energy Deficit Window causes localized NAD+ insufficiency that compromises the metabolically expensive restoration machinery, lowering $k_r$ at the very window where $k_g(t)$ peaks --- driving the ratchet toward progressive cementing of curvature.

\subsection{Molecular Grounding: The AlphaFold Anisotropy Gap (Exploratory)}

While the primary failure mode is governed by macroscopic delay dynamics, we also explore the molecular architecture of the mechanosensory apparatus. AlphaFold v6 structural analysis of 23 key spinal proteins reveals that demand-side mechanosensory proteins (e.g., Vimentin, PIEZO2) exhibit significantly higher structural anisotropy (mean $3.08 \pm 1.44$) than supply-side metabolic regulators (mean $1.79 \pm 0.56$).

We tentatively propose a \textbf{Demand-Supply Anisotropy Gap} ($p = 0.011$ in our dataset): the proteins required to \emph{sense} mechanical geometry are structurally extended and thermodynamically expensive to maintain, while their metabolic regulators are compact and often intrinsically disordered. We emphasize that this observation is hypothesis-generating. Current AI structural prediction tools cannot directly infer mechanical properties under load~\cite{gomes2022protein, zonta2024sequence}, and confidence scores for extracellular matrix components are frequently low. Nonetheless, this structural asymmetry provides a compelling direction for future experimental validation of the cellular mechanisms underlying the Energy Deficit Window.
